% Source: https://artofproblemsolving.com/community/c3h2653608_collatz_conjecture
% Stable identity: aops-topic-2653608-q01
\textbf{(a)} From an odd $n$ the two steps give $3n+1$ and then $\dfrac{3n+1}{2}$, which is \emph{larger} than $n$. Halving once does not undo multiplying by $3$, and the halved number may be odd again, so the sequence can keep growing for a while. Nothing in the argument shows that the terms eventually decrease, let alone reach $1$. (Whether every starting number reaches $1$ is the Collatz conjecture, which is still unproved.)

\textbf{(b)} Starting from $7$:
\[7,\ 22,\ 11,\ 34,\ 17,\ 52,\ 26,\ 13,\ 40,\ 20,\ 10,\ 5,\ 16,\ 8,\ 4,\ 2,\ 1.\]
The sequence rises above $7$ several times, and its largest term is $\boxed{52}$.
